\chapter{Atualização do escopo gratuito e da acessibilidade}

A versão final da SkillSwap não cobra dos usuários. Cadastro, troca de habilidades, conversa, agendamento, aula e geração do comprovante de conclusão permanecem disponíveis sem pagamento obrigatório. As rotas relacionadas a carteira, SwapCoins, premium e Stripe foram retiradas da navegação e bloqueadas na API. O código legado dessas integrações é mantido somente para rastreabilidade da evolução e não representa uma funcionalidade disponível.

\section{Preferências escolhidas pela pessoa}

Após o cadastro com Firebase, a plataforma apresenta uma etapa opcional de configuração. A pessoa pode escolher contraste intenso, tamanho do texto, canal preferido de comunicação, legendas quando disponíveis, descrição de imagens, VLibras, linguagem simplificada e operação por teclado. Essas respostas são preferências de uso e não servem para diagnosticar, classificar ou expor uma deficiência.

As cores, formas e marcadores visuais são apenas temas de interface. Não existe uma regra que associe automaticamente uma cor ou uma forma a autismo, daltonismo, cegueira, deficiência motora ou ausência de deficiência. Essa decisão preserva a autonomia da pessoa e evita que uma característica sensível seja inferida de maneira incorreta.

\section{Contraste e comunicação equivalente}

O contraste intenso utiliza fundo escuro, texto claro, foco visível, bordas destacadas e informações textuais junto aos elementos visuais. O vermelho não é usado como único indicador de erro. O sistema mantém texto, teclado e mensagens de estado como caminhos alternativos quando áudio, vídeo, câmera, microfone ou conexão externa não estiverem disponíveis.

VLibras e legendas são recursos complementares. A presença de um widget ou de uma opção de legenda não é apresentada como certificação WCAG. A avaliação completa depende de testes com teclado, leitores de tela, zoom, tecnologias assistivas e participantes voluntários.

\section{Conclusão da aula e comprovante}

A conclusão da aula depende da confirmação do aluno e do professor. Depois das duas confirmações, o aluno pode gerar uma página acessível de certificado contendo seu nome, a aula, o professor, a duração, a data e um identificador. O arquivo pode ser impresso ou salvo como PDF pelo navegador.

O comprovante registra participação na atividade da SkillSwap. Ele não equivale a diploma, certificado profissional ou validação externa de competência.
